\chapter{Introdução}
\label{cap:introducao}
A ascensão da Inteligência Artificial (IA) e, proeminentemente, dos algoritmos de Aprendizado de Máquina (Machine Learning), inaugurou uma nova era na governança tecnológica e na gestão de sistemas complexos. No âmbito da administração pública e da gestão educacional, a transição de processos puramente burocráticos para decisões orientadas a dados — impulsionada pela Mineração de Dados Educacionais (MDE) — tem sido amplamente adotada como estratégia para otimizar a alocação de recursos, prever taxas de evasão e estimar o desempenho acadêmico de estudantes. Essa engenharia preditiva baseia-se na premissa de que o processamento massivo de históricos populacionais permite antecipar vulnerabilidades e subsidiar formulações de políticas públicas baseadas em evidências.

Contudo, a sofisticação matemática desses modelos preditivos impõe um desafio crítico: o fenômeno da "caixa-preta" (black-box problem). Algoritmos de alta performance interpretativa frequentemente operam sob uma opacidade intrínseca, em que os caminhos lógicos que conduzem a uma determinada classificação ou predição são incompreensíveis para observadores humanos. Conforme aponta a literatura clássica de sistemas preditivos, a ausência de transparência e interpretabilidade cria um abismo de responsabilidade, impedindo que gestores, educadores e os próprios cidadãos afetados compreendam os critérios determinantes por trás de uma decisão automatizada.

Essa opacidade torna-se alarmante quando se constata que os modelos de aprendizado de máquina não são agentes neutros; eles herdam, decodificam e amplificam os vieses históricos presentes nos dados utilizados para o seu treinamento. Quando aplicados a bases de dados que refletem profundas desigualdades estruturais, os algoritmos tendem a converter disparidades socioeconômicas, de gênero e de raça em regras matemáticas de exclusão, penalizando sistematicamente grupos historicamente vulneráveis sob o manto de uma falsa objetividade técnica.  

Diante desse cenário, a Inteligência Artificial Explicável (XAI) e a Justiça Algorítmica (Algorithmic Fairness) emergem não apenas como subcampos de pesquisa computacional, mas como imperativos éticos e legais. No contexto brasileiro, regulamentações como a Lei Geral de Proteção de Dados (LGPD) já asseguram o direito do cidadão à revisão e à explicação de decisões tomadas unicamente por tratamentos automatizados. Torna-se imperativo, portanto, o desenvolvimento de ferramentas e artefatos de software focados na auditoria algorítmica contínua, capazes de abrir a caixa-preta dos modelos e quantificar objetivamente o viés preditivo em exames públicos de larga escala. Nesse panorama, este trabalho propõe o desenvolvimento de uma plataforma computacional de arquitetura desacoplada e a aplicação integrada do framework SHAP com métricas de impacto disparate como forma de garantir a transparência, a auditabilidade e a equidade socioeconômica na avaliação de modelos preditivos treinados com os microdados do Exame Nacional do Ensino Médio (ENEM).

\section{Motivação e Justificativa}
A crescente adoção de sistemas de tomadas de decisão automatizadas, impulsionadas por algoritmos de Aprendizado de Máquina (Machine Learning), tem reconfigurado a gestão de políticas públicas e a alocação de recursos estruturais na sociedade contemporânea. No campo da educação, essa tendência se consolida por meio da Mineração de Dados Educacionais, na qual algoritmos são rotineiramente treinados para prever o desempenho acadêmico de estudantes, antecipar taxas de evasão e subsidiar processos de seleção. No cenário brasileiro, o epicentro analítico dessas iniciativas reside nos microdados do Exame Nacional do Ensino Médio (ENEM). Por se tratar do principal mecanismo de estratificação e acesso ao ensino superior no país, os dados do exame são amplamente utilizados para alimentar modelos preditivos que visam simular o sucesso acadêmico ou direcionar políticas de assistência. É precisamente nessa transição que o impacto de predições enviesadas assume proporções críticas: a aparente neutralidade matemática desses classificadores frequentemente oculta e perpetua as profundas desigualdades socioeconômicas e históricas já consolidadas na base de dados de treinamento.

Sob a ótica social, a relevância deste trabalho fundamenta-se na necessidade de assegurar a justiça algorítmica em exames de larga escala. Modelos preditivos treinados com os microdados do INEP correm o risco iminente de converter disparidades socioeconômicas históricas, tais como abismos de renda familiar, assimetrias digitais e desigualdades de raça/cor, em critérios automatizados de penalização estatística. Auditar esses modelos não é apenas um exercício de validação de métricas de acurácia, mas um imperativo ético para evitar a criação de barreiras invisíveis ao desenvolvimento acadêmico e profissional de minorias representadas nos dados.

Do ponto de vista científico e tecnológico, a justificativa deste estudo repousa sobre o desafio do dilema entre a interpretabilidade e a performance dos modelos de inteligência artificial. Algoritmos complexos de alta performance operam frequentemente sob a dinâmica de "caixas-pretas". A ausência de mecanismos transparentes que discriminem os fatores determinantes para a classificação de um estudante impede o diagnóstico preciso de padrões discriminatórios. Ao mobilizar técnicas de Inteligência Artificial Explicável (Explainable AI - XAI), especificamente por meio do cálculo de valores de Shapley via framework SHAP, este trabalho contribui para o avanço da Engenharia de Dados aplicada à auditoria computacional, demonstrando que é possível extrair explicações locais e globais robustas sem sacrificar a capacidade preditiva do artefato tecnológico.

Ademais, a pesquisa encontra forte eco no cenário regulatório brasileiro. A Lei Geral de Proteção de Dados (LGPD - Lei nº 13.709/2018), em seu Artigo 20, assegura ao titular dos dados o direito de solicitar a revisão de decisões tomadas unicamente com base em tratamento automatizado de dados pessoais que afetem seus interesses \cite{brasil2018lgpd}. Paralelamente, os debates legislativos sobre o Marco Legal da Inteligência Artificial no Brasil propõem diretrizes de governança que incluem avaliações de impacto algorítmico e mecanismos de transparência para sistemas classificados como de alto risco \cite{brasil2023pl2338}.

Por fim, no âmbito da Engenharia de Sistemas de Informação, a justificativa consolida-se na entrega de um artefato de software prático, escalável e de código aberto. Ciente das severas limitações computacionais de ambientes de nuvem públicos e gratuitos (free tier), este trabalho justifica-se pelo desenvolvimento de uma arquitetura otimizada e desacoplada, que separa o pipeline pesado de processamento analítico em lote (batch processing) da camada de visualização interativa. O trabalho deixa de ser uma mera análise estatística estática para se transformar em um ecossistema funcional de auditoria, replicável para outras bases de dados públicas governamentais, fornecendo a órgãos de controle, pesquisadores e formuladores de políticas públicas uma ferramenta tangível para a promoção de uma IA ética e transparente no ecossistema brasileiro.

\section{Objetivos}

\subsection{Objetivo Geral}
O objetivo geral deste trabalho é projetar, desenvolver e validar uma ferramenta de software de auditoria algorítmica capaz de detectar vieses analíticos em modelos preditivos de desempenho educacional, utilizando métricas quantitativas de justiça (Fairness) e provendo explicações centradas em Inteligência Artificial Explicável (XAI) no contexto de exames públicos em larga escala.

\subsection{Objetivos Específicos}
\begin{itemize}
    \item Construir um pipeline de dados (ETL) automatizado e otimizado para a ingestão, tratamento e persistência em formato Parquet dos Microdados do ENEM;
    \item Treinar, otimizar e homologar um classificador binário baseado no algoritmo XGBoost para servir de objeto alvo para a auditoria de viés;
    \item Implementar um motor computacional de Fairness para o cálculo automatizado das métricas de Paridade Demográfica e Impacto Disparate;
    \item Integrar uma camada de explicabilidade post-hoc baseada no framework SHAP para rastrear e quantificar a contribuição das variáveis socioeconômicas na decisão do classificador;
    \item Desenvolver uma interface gráfica interativa baseada no framework Streamlit, orientada de forma otimizada para o consumo assíncrono dos resultados analíticos;
    \item Validar a eficácia e a sensibilidade da ferramenta por meio da metodologia de injeção intencional de vieses sintéticos em cenários controlados.
\end{itemize}

\section{Estrutura do Trabalho}
Além desta introdução, o trabalho está organizado em mais seis capítulos:
\begin{itemize}
    \item \textbf{Capítulo~\ref{cap:referencial} --- Referencial Teórico:} apresenta os fundamentos de Mineração de Dados Educacionais, o problema da opacidade algorítmica, o algoritmo XGBoost, as definições formais de Justiça Algorítmica (Paridade Demográfica, Impacto Disparate, \textit{Equal Opportunity}), a explicabilidade baseada em valores de Shapley e a \textit{Design Science Research};
    \item \textbf{Capítulo~\ref{cap:trabalhos_relacionados} --- Trabalhos Relacionados:} discute estudos preditivos sobre o ENEM, ferramentas de \textit{fairness}, explicabilidade e o panorama regulatório brasileiro, e delimita a lacuna de pesquisa;
    \item \textbf{Capítulo~\ref{cap:metodologia} --- Metodologia:} detalha a aplicação da DSR, o tratamento dos dados, o modelo-alvo, o protocolo de auditoria restrito ao conjunto de teste e o experimento de validação por injeção sintética de viés;
    \item \textbf{Capítulo~\ref{cap:desenvolvimento} --- Desenvolvimento do Artefato:} descreve os requisitos, a arquitetura desacoplada, o pipeline em lote, o contrato de dados, as decisões de desempenho, a interface interativa e o empacotamento em Docker;
    \item \textbf{Capítulo~\ref{cap:resultados} --- Resultados e Discussão:} apresenta a auditoria das edições de 2019 a 2021 do ENEM (desempenho, equidade por raça/cor e renda, amplificação do viés, explicabilidade) e a validação do artefato;
    \item \textbf{Capítulo~\ref{cap:conclusao} --- Conclusão:} sintetiza as contribuições, as limitações e os trabalhos futuros.
\end{itemize}